\documentclass[11pt,a4paper]{article}
\usepackage[margin=2.3cm]{geometry}
\usepackage{graphicx}
\usepackage{amsmath}
\usepackage{amssymb}
\usepackage{booktabs}
\usepackage[hidelinks]{hyperref}
\usepackage{microtype}
\title{A neutral parameter sweep: no effective size reproduces the empirical
differentiation at the ancestry-tract-matched time}
\author{}
\date{Module E working note --- \today}

\begin{document}
\maketitle
\vspace{-2em}

\begin{abstract}\noindent
We ran the recombination-matched, haplodiploid neutral null as a parameter sweep in
carrying capacity $K$ and founding size, sampling each of 20 independent demes every 25
ticks to 250. Two results follow. (1) \textbf{Effective size} scales with carrying
capacity as $N_e \approx 0.31\,K$ (haplodiploidy and the age-structured life table pull
$N_e$ well below census $K$); founding size (100--200) barely affects $N_e$. (2) Using
within-deme adjacent-marker LD as an ancestry-tract clock, the empirical value is matched
by \emph{every} parameter combination at only $\sim$15--20 generations of recombination.
\textbf{At that tract-matched time the neutral among-population $F_{ST}$ is
$\sim$0.03--0.04, against the empirical 0.30 --- a 6--10$\times$ gap --- and no
combination, at any time out to 74 generations including the strongest drift, comes near
it.} Insufficient recombination is excluded, and no neutral effective size in this range
rescues the differentiation. This is the parameter-swept, direct form of the
$F_{ST}$--tract dissociation.
\end{abstract}

\section{The sweep}
Base configuration $K = 6250$, total founders $= 150$ (drawn from the diversified
1200/520 pool, split aquilonia-males : polyctena-queens to hold admixture proportion
$\alpha \approx 0.54$). One-at-a-time: $K \in \{3125, 6250, 12500\}$ and total founders
$\in \{100, 150, 200\}$, others at base --- five combinations, 20 demes each, one
replicate. Every deme runs to 250 ticks and is sampled every 25. Markers are the
\textbf{9{,}902 high-DiagnosticIndex loci} (DI $>-25$, MAF $\ge 0.15$, systematically
thinned to $\sim$15\,kb spacing), the ancestry-informative set that tract-length
estimation requires. \textbf{Ticks are not generations}: the life table gives a
generation time of $\approx 3.4$ ticks, so 250 ticks $\approx 74$ generations; all
generation axes below are ticks$/3.4$.

\section{Effective size (Fig.~\ref{fig:sweep}A, Table~\ref{tab:ne})}
Within-deme expected heterozygosity decays log-linearly with generation, giving $N_e$
from the slope ($H_e(g) = H_e(0)\,e^{-g/2N_e}$).

\begin{table}[h]\centering
\caption{Effective size from within-deme heterozygosity decay (generations $=$ tick$/3.4$).}
\label{tab:ne}
\begin{tabular}{lcccc}
\toprule
combination & $K$ & founders & $N_e$ & fit $R^2$ \\
\midrule
K3125\_f150  & 3125  & 150 & 1{,}195 & 0.99 \\
K6250\_f150  & 6250  & 150 & 1{,}950 & 0.99 \\
K12500\_f150 & 12500 & 150 & 3{,}917 & 0.93 \\
K6250\_f100  & 6250  & 100 & 1{,}980 & 0.99 \\
K6250\_f200  & 6250  & 200 & 2{,}105 & 0.99 \\
\bottomrule
\end{tabular}
\end{table}

$N_e \approx 0.31\,K$ across the $K$ sweep. Founding size (100--200) changes $N_e$
negligibly: the bottleneck sets the \emph{initial} differentiation, not the long-run
drift rate. So $K$ is the effective-size lever and founding size is a starting-condition
lever. Note that every simulated deme retains $H_e \approx 0.47$, whereas the empirical
value on the same markers is $H_e = 0.31$ (dashed) --- the empirical diagnostic loci are
far more sorted (diversity-depleted) than neutral drift produces.

\section{The tract-matched dissociation (Fig.~\ref{fig:sweep}B,C)}
Panel C uses within-deme adjacent-marker $r^2$ as an ancestry-tract clock: it is the
admixture LD that recombination erodes as tracts shorten. The empirical value (0.556) is
crossed by \emph{all five} combinations at $\sim$15--20 generations (ticks $\sim$50--75);
the curves overlie, so the tract clock is nearly parameter-independent and reads
$\sim$15--20 generations --- consistent with the $\sim$20-generation field age of the
hybrid populations.

Panel B is the differentiation at those same times. At the tract-matched
$\sim$15--20 generations the simulated among-deme $F_{ST}$ is $\sim$0.03--0.04, against
the empirical \textbf{0.30}. The gap is 6--10$\times$, and it does not close anywhere in
the sweep: even the strongest drift (K3125, $N_e = 1195$) at 74 generations reaches only
$F_{ST} = 0.055$, by which time its tracts (adjacent LD 0.47) have already fallen
\emph{below} empirical. For drift alone to reach $F_{ST} = 0.30$ in $\sim$20 generations
would require $N_e \approx 30$ ($K \approx 100$) --- implausible, and it would deplete
diversity genome-wide rather than concentrate differentiation at diagnostic loci.

\begin{figure}[h]\centering
\includegraphics[width=\textwidth]{../Figures/moduleE_sweep_summary.pdf}
\caption{Parameter sweep vs generations ($=$ ticks$/3.4$); empirical on the same 9{,}902
high-DI markers is dashed. \textbf{(A)} within-deme diversity decay (source of $N_e$);
all sims sit far above the empirical (heavily sorted) value. \textbf{(B)} among-deme
$F_{ST}$ --- every combination is $\sim$0.03--0.05, an order of magnitude below empirical
0.30. \textbf{(C)} within-deme adjacent-marker LD (tract clock) --- the empirical value
is matched at $\sim$15--20 generations. Reading B at the C-matched time gives the
6--10$\times$ $F_{ST}$ deficit.}
\label{fig:sweep}
\end{figure}

\section{Conclusion}
Across a $K \times$ founding-size grid spanning $N_e \approx 1200$--3900, no neutral
combination reproduces the empirical among-population differentiation at diagnostic loci
at the time its ancestry tracts match empirical --- nor at any time in the 74-generation
window. The differentiation deficit is 6--10$\times$ and parameter-robust, and the
diversity is simultaneously far too high. This is the direct, parameter-swept statement
of the result the earlier analyses inferred: \textbf{the empirical high-$F_{ST}$,
tract-appropriate, diversity-depleted state at ancestry-informative loci is not a neutral
outcome at any effective size or age in this range; it is the signature of locus-specific
ancestry sorting.} Insufficient recombination, insufficient drift, and founding-bottleneck
strength are all excluded as explanations.

\paragraph{Preliminary in one respect.} Panel C uses within-deme adjacent-marker LD as a
\emph{proxy} for tract length. The colleagues' ancestry-tract caller will give the actual
tract-length distribution (in cM), pinning the matched generation more precisely and
allowing the full distribution shape to be matched rather than a single LD summary. The
$F_{ST}$--tract dissociation is already unambiguous and does not depend on the proxy.

\section*{Reproducibility}
Sweep harness and inputs: \texttt{formica\_neutral\_sim\_sweep\_handoff/} (README,
\texttt{run\_sweep.sh}, model, \texttt{markers.tsv}, founders). Output:
\texttt{moduleE\_slim/output\_track\_length/} (5 combinations $\times$ 20 demes $\times$
10 ticks). Analysis: \texttt{dev/R/moduleE\_analyze\_sweep.R} (Fig.~\ref{fig:sweep},
Table~\ref{tab:ne}). Generation time from the life table; $\sim$3.4 ticks/generation.

\end{document}
